\chapter{Conclusion}

For the same 11,719 RetailRocket purchasers, compact and rich features produce substantially different K Means and CLASSIX assignments. Matched-parameter comparisons also show disagreement, so the result does not arise only from selecting different cluster counts. High resampling agreement shows that each configuration repeats its own grouping, without implying agreement across feature spaces.

Representation also changes the explanation evidence. Rich ExKMC fidelity begins higher at four leaves, but compact fidelity exceeds it at sixteen leaves. Thus the relative accuracy of these rule approximations depends on the permitted complexity. The rich CLASSIX fit has more intermediate aggregation groups, which describes the detail of its construction record rather than measured human difficulty. One rich K Means cluster passes the size, coverage and distinctiveness audit; no compact K Means or CLASSIX cluster passes. Accurate rule approximation and repeatable clustering therefore do not by themselves establish operational addressability.

The sensitivity analyses qualify these findings under feature deletion, reduced dimension, numerical perturbation and alternative session windows. They do not show that every representation choice would preserve the same result. The projection argument explains a valid computational shortcut, while the local stability proposition gives sufficient conditions only for a fixed distance-merging fit. It does not certify the selected rich density-merging result.

Representation is part of the empirical clustering model rather than a neutral preprocessing choice. In this dataset, richer behavioural information changes geometry, assignments and explanation properties, without consistently improving clustering quality, simplifying explanations or producing deployment-ready groups. The study does not establish objectively correct customer types, human comprehension or realised business value. Those questions require evidence beyond the recorded event log.
